% Сборка:
% lualatex -halt-on-error -interaction=nonstopmode grassmann_rank_one.tex
\documentclass[11pt,a4paper]{article}

\usepackage{fontspec}
\setmainfont{Noto Serif}
\usepackage[russian]{babel}
\usepackage[margin=20mm]{geometry}
\usepackage{amsmath,amssymb,amsthm,mathtools,microtype}
\usepackage{enumitem}

\setlength{\parindent}{0pt}
\setlength{\parskip}{4pt}
\setlength{\emergencystretch}{2em}

\newcommand{\R}{\mathbb R}
\newcommand{\T}{^{\mathsf T}}
\newcommand{\F}{_{\mathrm F}}
\newcommand{\Id}{\mathrm{Id}}
\newcommand{\Gr}{\mathrm{Gr}}
\newcommand{\St}{\mathrm{St}}
\newcommand{\Exp}{\operatorname{Exp}}
\newcommand{\grad}{\operatorname{grad}}
\newcommand{\spanop}{\operatorname{span}}
\newcommand{\tr}{\operatorname{tr}}
\DeclareMathOperator{\rank}{rank}
\DeclareMathOperator{\argmin}{argmin}
\DeclareMathOperator{\argmax}{argmax}

\newtheorem{proposition}{Утверждение}
\newtheorem{remark}{Замечание}

\begin{document}

\begin{center}
{\Large\bfseries Жадная оптимизация на Grassmann-многообразии}\\[4pt]
{\large Последовательность rank-one геодезических поворотов для multi-index задачи}
\end{center}

\section{Исходная multi-index задача}

Пусть
\[
    P\in\R^{m\times d},
    \qquad
    PP\T=\Id_m,
    \qquad
    m\ll d,
\]
задаёт текущую оценку EDR-подпространства. Удобнее перейти к
\[
    Y=P\T\in\R^{d\times m},
    \qquad
    Y\T Y=\Id_m.
\]
Матрицы \(Y\) и \(YR\), где \(R\in O(m)\), задают одно и то же
\(m\)-мерное подпространство. Поэтому естественным параметром служит
\[
    [Y]\in \Gr(m,d).
\]

Для локальных статистик
\[
    I_j\in\R^p,
    \qquad
    U_j\in\R^{p\times d},
    \qquad
    N_j>0,
    \qquad
    j=1,\ldots,J,
\]
исходный квадратичный функционал имеет вид
\[
    L(Y,\{g_j\})
    =
    \frac14
    \sum_{j=1}^J
    N_j
    \left\|
        I_j-U_jYg_j
    \right\|_2^2,
    \qquad
    g_j\in\R^m.
\]
Здесь \(g_j\) являются координатами локального градиента в выбранном
базисе \(Y\).

Если заменить \(Y\) на \(YR\), то необходимо одновременно заменить
\[
    g_j\mapsto R\T g_j.
\]
Поэтому функционал с фиксированными \(g_j\) не является функцией только
от подпространства. Для корректной Grassmann-постановки локальные
коэффициенты следует исключить.

\section{Профилированный функционал на \(\Gr(m,d)\)}

Определим
\[
    g_j^\star(Y)
    =
    \argmin_{g\in\R^m}
    \left\{
        \|I_j-U_jYg\|_2^2+\tau\|g\|_2^2
    \right\},
    \qquad
    \tau\ge 0.
\]
При невырожденности соответствующей системы, либо при \(\tau>0\),
\[
    g_j^\star(Y)
    =
    A_j(Y)^{-1}b_j(Y),
\]
где
\[
    A_j(Y)
    =
    Y\T U_j\T U_jY+\tau\Id_m,
    \qquad
    b_j(Y)
    =
    Y\T U_j\T I_j.
\]

Введём профилированный функционал
\[
    \Phi_{\rm data}(Y)
    =
    \frac14
    \sum_{j=1}^J
    N_j
    \min_{g\in\R^m}
    \left\{
        \|I_j-U_jYg\|_2^2+\tau\|g\|_2^2
    \right\}.
\]
Эквивалентно,
\[
    \Phi_{\rm data}(Y)
    =
    \frac14
    \sum_{j=1}^J
    N_j
    \left[
        \|I_j\|_2^2
        -
        b_j(Y)\T A_j(Y)^{-1}b_j(Y)
    \right].
\]

\begin{proposition}[Grassmann-инвариантность]
Для любого \(R\in O(m)\)
\[
    \Phi_{\rm data}(YR)=\Phi_{\rm data}(Y).
\]
Следовательно, \(\Phi_{\rm data}\) корректно задаёт функцию на
\(\Gr(m,d)\).
\end{proposition}

\begin{proof}
Замена \(Y\mapsto YR\) даёт взаимно однозначную замену координат
\(g\mapsto R\T g\), причём
\[
    U_jYR(R\T g)=U_jYg
\]
и \(\|R\T g\|_2=\|g\|_2\). Поэтому значение внутреннего минимума
не меняется.
\end{proof}

Если требуется удерживать новое подпространство около текущего
\(Y_0=P\T\), штраф Фробениуса между конкретными базисами
\(\|Y-Y_0\|\F^2\) использовать нельзя: он зависит от выбора базиса.
Вместо него можно использовать квадрат chordal-расстояния
\[
    d_{\rm ch}^2(Y,Y_0)
    =
    \frac12
    \|YY\T-Y_0Y_0\T\|\F^2
    =
    m-\|Y_0\T Y\|\F^2.
\]
Полный функционал:
\[
    \boxed{
    \Phi(Y)
    =
    \Phi_{\rm data}(Y)
    +
    \lambda
    \left(
        m-\|Y_0\T Y\|\F^2
    \right).
    }
\]

\section{Локальный Taylor-аналог на Grassmann-многообразии}

Касательное пространство в точке \(Y\) удобно представлять как
\[
    T_Y\Gr(m,d)
    =
    \left\{
        Z\in\R^{d\times m}:
        Y\T Z=0
    \right\}.
\]
Для малого \(Z\in T_Y\Gr(m,d)\) экспоненциальное отображение задаёт
точку
\[
    Y_+=\Exp_Y(Z)\in\Gr(m,d).
\]

Локальное разложение геометрии имеет вид
\[
    \boxed{
    \Exp_Y(Z)
    =
    Y+Z-\frac12YZ\T Z
    +O(\|Z\|\F^3).
    }
\]
Первый порядок \(Y+Z\) является касательной линейной аппроксимацией,
а член
\[
    -\frac12YZ\T Z
\]
является первой поправкой на кривизну Grassmann-многообразия.

Для гладкого функционала:
\[
    \Phi(\Exp_Y(Z))
    =
    \Phi(Y)
    +
    \langle G,Z\rangle_F
    +
    \frac12
    \operatorname{Hess}\Phi(Y)[Z,Z]
    +
    O(\|Z\|\F^3),
\]
где
\[
    G=\grad\Phi(Y)\in T_Y\Gr(m,d).
\]
Это прямой аналог Taylor-разложения функции в евклидовом пространстве.

\section{Riemannian-градиент}

Пусть
\[
    r_j(Y)
    =
    I_j-U_jYg_j^\star(Y).
\]
По теореме об огибающей евклидов градиент data-части равен
\[
    \nabla \Phi_{\rm data}(Y)
    =
    -\frac12
    \sum_{j=1}^J
    N_j
    U_j\T r_j(Y)
    \bigl(g_j^\star(Y)\bigr)\T.
\]
Для штрафа
\[
    \Phi_{\rm prox}(Y)
    =
    \lambda
    \left(
        m-\|Y_0\T Y\|\F^2
    \right)
\]
имеем
\[
    \nabla\Phi_{\rm prox}(Y)
    =
    -2\lambda Y_0Y_0\T Y.
\]
Следовательно,
\[
    \nabla\Phi(Y)
    =
    -\frac12
    \sum_{j=1}^J
    N_j
    U_j\T r_j
    (g_j^\star)\T
    -
    2\lambda Y_0Y_0\T Y.
\]
Grassmann-градиент получается проекцией на касательное пространство:
\[
    \boxed{
    G
    =
    \grad\Phi(Y)
    =
    (\Id_d-YY\T)\nabla\Phi(Y).
    }
\]

Полную матрицу \(G\) строить необязательно. Для \(a\in\R^m\)
\[
    Ga
    =
    -\frac12
    (\Id_d-YY\T)
    \sum_j
    N_j
    ((g_j^\star)\T a)\,
    U_j\T r_j
    -
    2\lambda
    (\Id_d-YY\T)
    Y_0Y_0\T Ya,
\]
а для \(v\in\R^d\)
\[
    G\T v
    =
    -\frac12
    \sum_j
    N_j
    g_j^\star
    \bigl(r_j\T U_j(\Id_d-YY\T)v\bigr)
    -
    2\lambda
    Y\T Y_0Y_0\T
    (\Id_d-YY\T)v.
\]
Поэтому поиск ведущего rank-one направления можно реализовать
матрично-свободно.

\section{Rank-one касательное направление}

Рассмотрим
\[
    Z=va\T,
\]
где
\[
    v\in\R^d,
    \qquad
    a\in\R^m,
    \qquad
    \|v\|_2=\|a\|_2=1,
    \qquad
    Y\T v=0.
\]
Тогда
\[
    Y\T Z
    =
    Y\T va\T=0,
\]
то есть
\[
    Z\in T_Y\Gr(m,d),
    \qquad
    \rank(Z)=1.
\]

Геометрический смысл \(a\): внутри текущего EDR-подпространства
выбирается единичное направление
\[
    y=Ya.
\]
Вектор \(v\) выбирается в ортогональном дополнении
\[
    v\perp\spanop(Y).
\]
Таким образом, \(Z=va\T\) задаёт изменение только в двумерной плоскости
\[
    \spanop\{Ya,v\}.
\]

\section{Точная экспонента rank-one шага}

\begin{proposition}[Rank-one Grassmann rotation]
Пусть
\[
    Y\T Y=\Id_m,
    \qquad
    Y\T v=0,
    \qquad
    \|v\|_2=\|a\|_2=1.
\]
Для
\[
    Z=va\T
\]
точный геодезический шаг длины \(\eta\) равен
\[
    \boxed{
    \Exp_Y(\eta va\T)
    =
    Y(\Id_m-aa\T)
    +
    \bigl(
        Ya\cos\eta+v\sin\eta
    \bigr)a\T.
    }
\]
\end{proposition}

\begin{proof}
Разложим базис \(Y\) на компоненту вдоль \(a\) и ортогональную к ней:
\[
    Y
    =
    Y(\Id_m-aa\T)
    +
    (Ya)a\T.
\]
Экспоненциальное отображение для rank-one касательного вектора
вращает только пару ортонормированных векторов
\[
    Ya,\qquad v
\]
в их двумерной плоскости:
\[
    Ya
    \mapsto
    Ya\cos\eta+v\sin\eta.
\]
Все направления \(Yb\), \(b\perp a\), остаются неизменными.
Полученная матрица сохраняет ортонормированность столбцов.
\end{proof}

Следовательно, один rank-one Grassmann-шаг является обычным
двумерным поворотом.

Его Taylor-разложение:
\[
\begin{aligned}
    \Exp_Y(\eta va\T)
    &=
    Y
    +
    \eta va\T
    -
    \frac{\eta^2}{2}Yaa\T
    -
    \frac{\eta^3}{6}va\T
    +
    O(\eta^4).
\end{aligned}
\]
Первый порядок совпадает с евклидовой rank-one поправкой
\[
    Y+\eta va\T,
\]
а следующие члены автоматически компенсируют кривизну и сохраняют
ортонормированность.

\section{Лучший rank-one шаг в первом порядке}

Пусть
\[
    G=\grad\Phi(Y).
\]
Ищем нормированный rank-one tangent vector, дающий максимальное
уменьшение линейной модели:
\[
    (v_\star,a_\star)
    \in
    \argmax_{\substack{
        \|v\|_2=\|a\|_2=1\\
        Y\T v=0
    }}
    \left|
        v\T G a
    \right|.
\]
Поскольку \(G\in T_Y\Gr(m,d)\), его ненулевые левые сингулярные
направления автоматически ортогональны \(Y\).

\begin{proposition}[Оптимальный rank-one tangent direction]
Если
\[
    G
    =
    U\Sigma V\T
\]
--- компактное SVD, то
\[
    \max_{\substack{
        \|v\|_2=\|a\|_2=1\\
        Y\T v=0
    }}
    |v\T Ga|
    =
    \sigma_1(G),
\]
и оптимальная пара задаётся ведущими сингулярными векторами \(G\).
Направление знака выбирается так, чтобы
\[
    \langle G,Z\rangle_F<0.
\]
\end{proposition}

Таким образом, полный Riemannian gradient step можно заменить
последовательностью жадных rank-one шагов.

Полное SVD \(G\) не обязательно. Ведущую пару можно искать
попеременными итерациями:
\[
    v
    \leftarrow
    -\frac{Ga}{\|Ga\|_2},
    \qquad
    a
    \leftarrow
    \frac{G\T v}{\|G\T v\|_2},
\]
с контролем знака и остановкой по стабилизации пары. Действия
\(Ga\) и \(G\T v\) вычисляются без формирования плотной
\(d\times d\) матрицы.

\section{Выбор угла поворота}

После выбора \(v,a\) остаётся одномерная задача
\[
    \eta_\star
    \in
    \argmin_{\eta\in[-\eta_{\max},\eta_{\max}]}
    \Phi\left(
        \Exp_Y(\eta va\T)
    \right).
\]
Это точный line search вдоль одной Grassmann-геодезической.

Более дешёвый вариант использует второй порядок:
\[
    \Phi(\Exp_Y(\eta Z))
    \approx
    \Phi(Y)
    +
    \eta\langle G,Z\rangle_F
    +
    \frac{\eta^2}{2}
    \operatorname{Hess}\Phi(Y)[Z,Z].
\]
Если
\[
    h_Z
    =
    \operatorname{Hess}\Phi(Y)[Z,Z]>0,
\]
то локально
\[
    \eta_\star
    \approx
    -\frac{\langle G,Z\rangle_F}{h_Z}.
\]
На практике можно использовать backtracking/Armijo без явного
вычисления Hessian.

\section{Последовательность простых поворотов}

Для текущей точки \(Y\) положим
\[
    y=Ya,
    \qquad
    \|y\|_2=1,
    \qquad
    y\perp v.
\]
Элементарный поворот можно записать как ортогональный оператор
\[
\begin{aligned}
    Q(y,v,\eta)
    &={}
    \Id_d
    +(\cos\eta-1)(yy\T+vv\T)\\
    &\quad
    +\sin\eta\,(vy\T-yv\T).
\end{aligned}
\]
Он действует как rotation на плоскости \(\spanop\{y,v\}\):
\[
    Qy=y\cos\eta+v\sin\eta,
    \qquad
    Qv=v\cos\eta-y\sin\eta,
\]
и тождественно действует на её ортогональном дополнении.

Для шага \(k\) определим
\[
    y_k=Y_ka_k,
    \qquad
    Q_k=Q(y_k,v_k,\eta_k).
\]
Тогда
\[
    \boxed{
    Y_{k+1}=Q_kY_k
    =
    Y_k(\Id_m-a_ka_k\T)
    +
    (Y_ka_k\cos\eta_k+v_k\sin\eta_k)a_k\T.
    }
\]
После \(K\) шагов получаем корректную композицию ортогональных
преобразований ambient-пространства:
\[
    \boxed{
    Y_K
    =
    Q_{K-1}Q_{K-2}\cdots Q_0Y_0.
    }
\]
Заметим, что \(Q_k\) строится уже из текущей точки \(Y_k\), поэтому
последовательность \(Q_k\) определяется адаптивно.

Каждый шаг сохраняет
\[
    Y_k\T Y_k=\Id_m,
\]
поэтому размерность подпространства всегда остаётся равной \(m\).
Никакого дополнения недостающих направлений после шага не требуется.

\section{Любой tangent step как семейство rank-one rotations}

Пусть произвольный касательный вектор имеет компактное SVD
\[
    Z
    =
    U\Sigma V\T
    =
    \sum_{\ell=1}^r
    \sigma_\ell u_\ell v_\ell\T,
    \qquad
    r=\rank(Z)\le m.
\]
Положим
\[
    y_\ell=Yv_\ell
\]
и введём кососимметричные генераторы
\[
    \Omega_\ell
    =
    u_\ell y_\ell\T-y_\ell u_\ell\T.
\]
Каждый \(\Omega_\ell\) действует только в двумерной плоскости
\[
    \spanop\{y_\ell,u_\ell\}.
\]
Из ортогональности сингулярных векторов следует, что эти плоскости
попарно ортогональны и
\[
    \Omega_\ell\Omega_s
    =
    \Omega_s\Omega_\ell,
    \qquad
    \ell\ne s.
\]

\begin{proposition}[Точное разложение Grassmann-геодезической]
Для
\[
    \Omega
    =
    \sum_{\ell=1}^r
    \sigma_\ell\Omega_\ell
\]
имеем
\[
    \Exp_Y(Z)
    =
    e^\Omega Y
    =
    \prod_{\ell=1}^r
    e^{\sigma_\ell\Omega_\ell}Y.
\]
Следовательно, любой один Grassmann-geodesic step раскладывается
\emph{точно}, а не только приближённо, в не более чем \(m\)
попарно коммутирующих двумерных rotations.
\end{proposition}

Это утверждение относится к разложению одного заранее заданного
касательного вектора \(Z\). Жадная схема отличается: после каждого
rank-one поворота градиент пересчитывается, поэтому следующие
\(v_k,a_k\) выбираются уже для новой точки \(Y_k\). Такая траектория
в общем случае не совпадает с одной геодезической
\(\Exp_{Y_0}(Z)\).

\section{Жадный rank-one Grassmann алгоритм}

Вход:
\[
    Y_0=P\T,
    \qquad
    Y_0\T Y_0=\Id_m.
\]

На шаге \(k\):

\begin{enumerate}[label=\arabic*.]
    \item Для текущего \(Y_k\) вычислить
    \[
        g_j^\star(Y_k),
        \qquad
        r_j(Y_k).
    \]

    \item Вычислять действия Riemannian gradient
    \[
        a\mapsto G_ka,
        \qquad
        v\mapsto G_k\T v
    \]
    без обязательного формирования всей матрицы \(G_k\).

    \item Найти приближённую ведущую rank-one пару
    \[
        (v_k,a_k)
        \approx
        \argmax_{\substack{
            \|v\|=\|a\|=1\\
            Y_k\T v=0
        }}
        |v\T G_ka|.
    \]
    Знак выбрать так, чтобы
    \[
        v_k\T G_ka_k<0.
    \]

    \item Если
    \[
        |v_k\T G_ka_k|
    \]
    меньше заданного допуска, остановиться.

    \item Найти угол \(\eta_k\) одномерным line search либо по
    локальной квадратичной модели.

    \item Выполнить точный rank-one rotation:
    \[
        \boxed{
        Y_{k+1}
        =
        Y_k(\Id_m-a_ka_k\T)
        +
        \bigl(
            Y_ka_k\cos\eta_k
            +
            v_k\sin\eta_k
        \bigr)a_k\T.
        }
    \]

    \item При необходимости из-за накопления round-off ошибок
    периодически выполнить thin QR. Это численная стабилизация,
    а не реконструкция потерянных направлений.
\end{enumerate}

\section{Связь с текущим SVD-решателем}

Текущий SVD-подход оптимизирует в ambient-пространстве матрицу
\[
    B\in\R^{m\times d}
\]
с ограничением
\[
    \rank(B)\le r<m
\]
либо малоранговую поправку \(B=P+\Delta\). В matrix-режиме
полученная матрица может иметь ранг меньше \(m\). Поэтому после
решения необходимо дополнять найденные направления до полного
\(m\)-мерного базиса.

В Grassmann-схеме состояние алгоритма всегда удовлетворяет
\[
    Y_k\in\St(d,m),
    \qquad
    [Y_k]\in\Gr(m,d).
\]
Rank-one ограничение относится не к рангу самого базиса, а к рангу
\emph{касательного изменения}
\[
    Z_k=v_ka_k\T.
\]
Следовательно,
\[
    \boxed{
    \rank(Y_k)=m
    \quad\text{на каждом шаге}.
    }
\]
Проблема восстановления \(m-r\) недостающих направлений отсутствует.

При этом Grassmann-схема не является механической заменой SVD в
текущем функционале с фиксированными \(g_j\). Для корректности требуется:

\begin{itemize}
    \item профилировать \(g_j\) или обновлять их согласованно с
    изменением базиса;
    \item заменить \(\|B-P\|\F^2\) на штраф между подпространствами;
    \item выполнять rank-one ограничение в касательном пространстве,
    а не на самой матрице базиса.
\end{itemize}

\section{Интерпретация как ``Taylor-семейства'' простых моделей}

Предлагаемая конструкция имеет два уровня.

Первый уровень является локальным:
\[
    \Gr(m,d)
    \longrightarrow
    T_Y\Gr(m,d),
\]
то есть сложное многообразие заменяется линейным касательным
пространством. Это аналог первого порядка Taylor-разложения.

Второй уровень:
\[
    Z
    =
    \sum_{\ell=1}^r
    \sigma_\ell v_\ell a_\ell\T,
\]
то есть касательный вектор раскладывается в rank-one компоненты.
Каждая компонента интегрируется на многообразии точным двумерным
поворотом.

Поэтому схема имеет вид
\[
    \boxed{
    \begin{gathered}
    \text{движение по }\Gr(m,d)
    \longrightarrow
    \text{касательная модель}\\
    \longrightarrow
    \text{rank-one направления}
    \longrightarrow
    \text{плоские rotations}
    \end{gathered}
    }
\]

В отличие от обычного Taylor-ряда, здесь rank-one rotations не только
дают локальную аппроксимацию. Для одного фиксированного tangent vector
его Grassmann-экспонента допускает точное разложение на ортогональные
двумерные rotations. При жадной оптимизации эти rotations выбираются
последовательно уже по новому градиенту.

\section{Что остаётся теоретически проверить}

Для применения к текущему multi-index алгоритму отдельно требуются
следующие результаты.

\begin{enumerate}[label=\arabic*.]
    \item Доказать корректность и регулярность профилированного
    функционала при возможной потере ранга локальных матриц
    \(U_jY\).

    \item Сравнить один жадный rank-one Grassmann-шаг с текущим
    rank-one SVD-шагом при малых изменениях \(Y\).

    \item Получить условие гарантированного убывания
    \[
        \Phi(Y_{k+1})\le \Phi(Y_k)
    \]
    для выбранного line search.

    \item Исследовать, достаточно ли последовательности ведущих
    rank-one направлений для достижения стационарной точки полного
    Grassmann-функционала.

    \item Сравнить вычислительную стоимость матрично-свободных действий
    \(Ga\), \(G\T v\) с текущими LSMR-решениями для \(v\).

    \item Проверить статистически, сохраняет ли профилированная
    Grassmann-постановка качество восстановления EDR-подпространства
    при тех же локальных статистиках \(I_j,U_j,N_j\).
\end{enumerate}

\end{document}
